\appendix
\section{Proof of the marginal coverage guarantee}
\label{app:proof}

Write $\tilde s(x,y)=s(x,y)-\tilde\delta_y$ for the corrected score, and recall from
\cref{sec:splits} that $\tilde\delta_\cdot$ is built from $\varphi(\cdot)$ and from the rows
$\mathcal{I}_{\mathrm{fit}}$, while $c$ is computed from the disjoint rows
$\mathcal{I}_{\mathrm{rec}}$.

\emph{Step 1: $\tilde s$ is a fixed function.} $\varphi$ is a function of the classifier's
weights, and $g_\theta$, $\lambda$ and $n_{\mathrm{cal}}$ are functions of
$(X_i,Y_i)_{i\in\mathcal{I}_{\mathrm{fit}}}$. None of them reads
$(X_i,Y_i)_{i\in\mathcal{I}_{\mathrm{rec}}}$ or the test point. Conditionally on
$\mathcal{I}_{\mathrm{fit}}$, therefore, $\tilde s$ is a deterministic map
$\mathcal{X}\times\mathcal{Y}\to\mathbb{R}$, fixed before any of the rows that will be used to
calibrate $c$ are seen. All statements below are conditional on
$\mathcal{I}_{\mathrm{fit}}$; since the bound obtained is uniform in it, the unconditional
bound follows by averaging.

\emph{Step 2: exchangeability survives.} If $(Z_1,\dots,Z_m)$ is exchangeable and $f$ is a
fixed measurable map, then $(f(Z_1),\dots,f(Z_m))$ is exchangeable, because for any permutation
$\pi$ the vector $(f(Z_{\pi(1)}),\dots,f(Z_{\pi(m)}))$ is the image under the same product map
of a vector equal in distribution to $(Z_1,\dots,Z_m)$. Applying this with
$Z_i=(X_i,Y_i)$ for $i\in\mathcal{I}_{\mathrm{rec}}\cup\{n+1\}$ and $f=\tilde s$ gives that
$\{\tilde s(X_i,Y_i)\}_{i\in\mathcal{I}_{\mathrm{rec}}}\cup\{\tilde s(X_{n+1},Y_{n+1})\}$ is
exchangeable.

\emph{Step 3: the split conformal bound.} For exchangeable real random variables
$V_1,\dots,V_{n_{\mathrm{rec}}},V_{n+1}$ and $c$ the empirical quantile of
$V_1,\dots,V_{n_{\mathrm{rec}}}$ at level
$\lceil (n_{\mathrm{rec}}+1)(1-\alpha)\rceil/n_{\mathrm{rec}}$, the standard argument
\cite{vovk2005algorithmic,angelopoulos2023gentle} gives
$\mathbb{P}(V_{n+1}\le c)\ge 1-\alpha$. Take $V_i=\tilde s(X_i,Y_i)$.

\emph{Step 4: the event is set membership.} By \cref{eq:set},
$y\in\mathcal{C}(x)\iff s(x,y)\le\hat q+\tilde\delta_y+c\iff\tilde s(x,y)\le \hat q+c$, so
absorbing the constant $\hat q$ into the threshold, $\{Y_{n+1}\in\mathcal{C}(X_{n+1})\}$ is
exactly $\{\tilde s(X_{n+1},Y_{n+1})\le \hat q + c\}$ and Step 3 applies. $\square$

\paragraph{What the proof does not give.} Only one scalar is calibrated on
$\mathcal{I}_{\mathrm{rec}}$, so $\tilde\delta_y$ may be arbitrarily complex without affecting
the argument; but for the same reason the conclusion is a statement about the marginal
distribution alone. Nothing here bounds
$\mathbb{P}(Y_{n+1}\in\mathcal{C}(X_{n+1})\mid Y_{n+1}=y)$ for any particular $y$, and for a
class with $n_y=0$ no such bound is available to any procedure. The disjointness in Step 1 is
also load-bearing: an earlier version of our code selected $\lambda$ over all classes including
the held-out ones, which does not merely bias an estimate but removes the premise, and it moved
the reported effect from $0.3$ to $0.5$ on the diagnostic we caught it with.

\section{Pre-registration and the fallback policy}
\label{app:prereg}

Two documents were frozen before any experiment in this paper was run, and both are in the
released repository. We summarise what they fix, because several of the choices would otherwise
look like decisions made after seeing results.

\paragraph{Measurability thresholds.} A survey of the score dumps, not of our memory, recorded
the median number of evaluation examples per class: $1168$ on ImageNet, $100$ on CIFAR-100,
$3$ on Pl@ntNet-300K and $2$ on iNaturalist-2018. With two examples a class's coverage can only
be $0$, $\tfrac12$ or $1$, so a minimum over $8142$ such classes is almost surely $0$ for every
procedure including the oracle, and differences between methods in that column carry no
information. The rule fixed in advance is therefore a median of at least $30$ evaluation
examples per class for per-class statistics (regime~A) and at least $200$ pooled evaluation rows
per prevalence bin otherwise (regime~B). Below the threshold, per-class statistics are still
computed but are written under a \texttt{withheld\_unmeasurable} key and may not enter a table
or decide a verdict. The same rule was later applied to $\lambda$ \emph{selection} as well as to
reporting, which was an omission we found on 2026-08-15 and fixed.

\paragraph{Behaviour of each baseline at $n_y=0$.} What each competitor emits for a class with
no calibration examples was written down before its released code was read: an infinite
threshold for classwise CP and for RC3P, a fall-back to the marginal threshold for Clustered CP
through its null cluster, and a kernel-weighted quantile over labelled classes for the method of
\citet{ding2026conformal}. All four predictions were then confirmed against the code. The
fall-back policy applied to an infinite threshold, namely substituting $\hat q$, was fixed at
the same time; it is what produces the zeros in \cref{tab:competitors}, and \cref{sec:comp}
says so in the text so that a zero is not read as a measurement.

\paragraph{Changes made after seeing results.} Three, all recorded. The oracle ceiling was
redefined from the unshrunk to the shrunk per-class quantile after PCC exceeded the unshrunk
version, for the reason given in \cref{app:oracle}. The calibration fraction was raised from
$0.50$ to $0.70$ after two points of the calibration-depth sweep turned out not to be binding on
a thin slice. The $\lambda$ ablation was extended from five splits to ten after it missed
Holm-corrected significance at five ($p=0.0094$ against a threshold of $0.0083$); it passes at
ten, and we report the number of splits for every table.

\section{The oracle ceiling, shrunk and unshrunk}
\label{app:oracle}

The ceiling reported in the main text is the \emph{best shrinkage} of a perfect per-class
offset: we form $\delta_y$ from the test labels and then select $\lambda$ for it exactly as
the method does. An earlier version used the unshrunk per-class quantiles, and PCC exceeded
it --- which is how we learned that the unshrunk version is not a bound. It is $\lambda=1$
applied to a perfect offset, and $\lambda=1$ is harmful for worst-class coverage
(\cref{tab:ablation}).

Across all $228$ configurations the two definitions agree in
$194$ cases, because the oracle's own shrinkage selects
$\lambda_{\text{oracle}}=1$ there; on the primary ImageNet arm they coincide at $+0.140$, so
the correction does not move that headline number. Where they differ, they differ by a lot,
and always in the same direction. \Cref{tab:oracle} lists the five largest gaps: the unshrunk
``ceiling'' falls as low as $-0.657$, far \emph{below} the marginal threshold it is supposed
to bound from above.

\begin{table}[h]
\centering
\small
\setlength{\tabcolsep}{4pt}
\begin{tabular}{lccc}
\toprule
Configuration & Shrunk & Unshrunk & $\lambda_{\text{oracle}}$ \\
\midrule
C\_comp\_competitors\_s1 & $+0.0487$ & $-0.6571$ & $0.10$ \\
H\_heldout\_ho0p5\_s2 & $+0.1928$ & $-0.4605$ & $0.70$ \\
C\_comp\_competitors\_s2 & $+0.1569$ & $-0.4902$ & $0.70$ \\
H\_heldout\_ho0p7\_s2 & $+0.0667$ & $-0.5733$ & $0.70$ \\
H\_heldout\_ho0p9\_s2 & $+0.1579$ & $-0.4605$ & $0.70$ \\
\bottomrule
\end{tabular}
\caption{The five configurations where the two ceiling definitions differ most. A quantity
that can sit $0.65$ below the baseline it is meant to bound is not a ceiling; the shrunk
version cannot be exceeded by PCC by construction, since PCC optimises the same $\lambda$
over a worse $\delta$.}
\label{tab:oracle}
\end{table}

\section{Why recalibration is invisible at matched size}
\label{app:invisible}

Let $t_y=\hat q+\tilde\delta_y+c$ be the thresholds a variant emits and let $T$ be the target
average set size. The matched comparison reports each variant at $t_y+\sigma$, where $\sigma$
solves $\frac1N\sum_i|\{y:s(X_i,y)\le t_y+\sigma\}|=T$. Adding a constant $a$ to every
threshold replaces $t_y$ by $t_y+a$ and hence $\sigma$ by $\sigma-a$, leaving
$t_y+a+(\sigma-a)=t_y+\sigma$ unchanged, so the matched thresholds are identical. Marginal recalibration adds exactly such a
constant, and so does the $n^\star$ rule when it substitutes one global value, which is why both
ablations must read $0.0000$ in \cref{tab:ablation}. The zero is a property of the comparison and
carries no information about the component.

\Cref{tab:invisible} gives the view in which the component is visible, namely the thresholds as
emitted, over ten ImageNet splits. Recalibration buys $0.0026$ of marginal coverage for $0.31$
of average set size. Two further readings matter. First, the raw \textsc{Pcc} thresholds are far
wider than the marginal ones, $2.31$ against $1.21$, so this view cannot be used to compare
methods and exists only to show where marginal validity lives; the size matching in the main
tables removes exactly this difference. Second, \textsc{Pcc} lands at $0.8989$ against a target
of $0.900$ while the marginal threshold lands at $0.9013$ on the same rows. The deficit of
$0.0024$ is about thirty times the finite-sample slack $1/(n+1)$ and is the measured cost of
letting $g_\theta$, $\lambda$ and $c$ share calibration rows, as discussed in
\cref{sec:guarantee}.

\begin{table}[h]
\centering
\small
\setlength{\tabcolsep}{4pt}
\begin{tabular}{lcc}
\toprule
Thresholds as emitted & Marginal cov. & Set size \\
\midrule
Marginal threshold $\hat q$        & $0.9013$ & $1.21$ \\
\textsc{Pcc}, no recalibration     & $0.8963$ & $2.00$ \\
\textsc{Pcc}, with recalibration   & $0.8989$ & $2.31$ \\
\midrule
Target                             & $0.9000$ & --- \\
\bottomrule
\end{tabular}
\caption{The unmatched view, ten ImageNet splits, $\alpha=0.10$. Recalibration moves marginal
coverage towards the target and costs set size. The matched comparison used everywhere else
cancels this shift by construction.}
\label{tab:invisible}
\end{table}

\section{Additional backbones, endogenous and exogenous descriptors}
\label{app:backbones}

We hypothesised that the two additional backbones failed because each was given its own
classifier head as $\varphi$, making the descriptors endogenous to the model that produced the
scores. \Cref{tab:backbones} refutes it: with a fixed ResNet-50 head instead --- whose weights
correlate $0.001$ element-wise with each model's own head, so it is genuinely exogenous ---
every interval still contains zero.

The reason is visible in the last column. All six arms have $35$ evaluation examples per class,
because ImageNet's validation split provides only $50$ per class in total and some must go to
calibration. \Cref{fig:evaldepth} places $35$ below the depth at which any effect is
detectable, so these experiments are inconclusive rather than negative. Settling the question
needs an evaluation set larger than the validation split.

\begin{table}[h]
\centering
\small
\setlength{\tabcolsep}{4pt}
\begin{tabular}{llccc}
\toprule
Backbone & $\varphi$ head & $\Delta$ worst & 95\% CI & Eval/class \\
\midrule
ConvNeXt-T \cite{liu2022convnet} & own & $+0.0000$ & $[+0.000,+0.000]$ & $35$ \\
ConvNeXt-T \cite{liu2022convnet} & fixed & $+0.0114$ & $[-0.017,+0.040]$ & $35$ \\
ResNet-50 \cite{he2016deep} & own & $+0.0171$ & $[-0.005,+0.040]$ & $35$ \\
ResNet-50 \cite{he2016deep} & fixed & $+0.0057$ & $[-0.005,+0.017]$ & $35$ \\
ViT-B/16 \cite{dosovitskiy2021image} & own & $-0.0286$ & $[-0.072,+0.015]$ & $35$ \\
ViT-B/16 \cite{dosovitskiy2021image} & fixed & $-0.0114$ & $[-0.066,+0.043]$ & $35$ \\
\bottomrule
\end{tabular}
\caption{Both descriptor arms on three backbones, five splits. ``own'' uses each model's own
classifier head, ``fixed'' the ResNet-50 head also used for ImageNet. Top-1 accuracies of all
three matched published values to four decimal places, which we used as a correctness check on
the weight extraction.}
\label{tab:backbones}
\end{table}

\section{Long-tailed datasets under a measurability threshold}
\label{app:longtail}

\begin{table}[h]
\centering
\small
\setlength{\tabcolsep}{3pt}
\begin{tabular}{llccc}
\toprule
Dataset & Thr. & Classes kept & $\Delta$ worst-class & $\lambda$ \\
\midrule
Pl@ntNet & $\ge 35$ & $151/1081$ & $+0.020\ [-0.011,+0.052]$ & $0.22$ \\
Pl@ntNet & $\ge 75$ & $\ \ 98/1081$ & $\mathbf{+0.019}\ [+0.010,+0.028]$ & $0.22$ \\
iNat-2018 & $\ge 35$ & $237/8142$ & $+0.000$ & $0.00$ \\
iNat-2018 & $\ge 75$ & $\ \ 75/8142$ & $+0.000$ & $0.00$ \\
\bottomrule
\end{tabular}
\caption{Restricting evaluation to classes whose slice can support a per-class estimate. Five
splits. Pl@ntNet recovers a significant improvement; iNaturalist does not, and there
$\lambda=0$ means the method declines to act rather than acting and failing.}
\label{tab:longtail}
\end{table}

The threshold changes the label space, so the row counts matter. On Pl@ntNet-300K the
$\ge 75$ rule keeps $98$ of $1081$ classes, of which $34$ are held out, and the survivors carry
a median of $116$ calibration and $198$ evaluation examples each. At $\ge 35$ it keeps $151$
classes with medians of $69$ and $116$. The improvement is significant at the stricter threshold
and not at the looser one, which is the ordering \cref{fig:evaldepth} predicts, since $35$ sits
at the edge of measurability while $75$ is clear of it.

On iNaturalist-2018 the survivors are not starved either: at $\ge 35$ the $237$ surviving
classes hold a median of $40$ calibration and $56$ evaluation examples. \Cref{tab:inatcurve}
gives the shrinkage curve on its labelled classes and explains the zero. The curve decreases
monotonically, so $\hat\delta$ is not degenerate; it carries signal, and the signal is wrong.
Because the grid contains $\lambda=0$, the selection step returns the marginal threshold and the
method neither helps nor harms. The cross-validated error of $g_\theta$ is $0.127$ here against
$0.036$ on ImageNet, so the regression is roughly three and a half times worse in a label space
eight times larger, fitted in a descriptor space of the same dimension. We did not test whether
richer descriptors recover the signal and state that as the open question.

\begin{table}[h]
\centering
\small
\setlength{\tabcolsep}{4pt}
\begin{tabular}{lcccccc}
\toprule
$\lambda$ & $0$ & $0.1$ & $0.2$ & $0.3$ & $0.5$ & $1.0$ \\
\midrule
worst-class & $\mathbf{0.699}$ & $0.466$ & $0.408$ & $0.385$ & $0.350$ & $0.319$ \\
\bottomrule
\end{tabular}
\caption{Shrinkage curve on the labelled classes of iNaturalist-2018, mean over five splits, at
the $\ge 75$ threshold. Every positive $\lambda$ is worse than none, so the selection step
returns $\lambda=0$. A flat curve would have meant $\hat\delta$ was degenerate; a decreasing one
means it is informative and points the wrong way.}
\label{tab:inatcurve}
\end{table}

\section{ImageNet-C, per corruption}
\label{app:imagenetc}

\Cref{tab:imagenetc} reports all $15$ corruptions at severities $3$ and $5$, three splits
each. Calibration uses clean images and evaluation uses corrupted ones, so exchangeability
fails and no method retains a coverage guarantee here (\cref{sec:shift}).

The empty-set column is the direct evidence of that. It ranges from
$0.063$ (brightness, severity 3) to
$0.757$ (contrast, severity 5), and an
average set size below $1.0$ is only possible when a large share of test points receive
\emph{nothing}. A threshold calibrated on clean images does not give narrow sets on corrupted
ones; it abstains. Coverage and set size both merely look ``low'' in that regime, which is why
we added this column.

\begin{table}[h]
\centering
\small
\setlength{\tabcolsep}{3pt}
\begin{tabular}{llcccc}
\toprule
Corruption & Sev. & $\Delta$ worst & Coverage & Set size & Empty \\
\midrule
brightness & 3 & $+0.0200$ & $0.841$ & $1.345$ & $0.063$ \\
brightness & 5 & $+0.0067$ & $0.746$ & $1.362$ & $0.103$ \\
contrast & 3 & $-0.0133$ & $0.600$ & $1.297$ & $0.171$ \\
contrast & 5 & $+0.0000$ & $0.053$ & $0.317$ & $0.757$ \\
defocus\_blur & 3 & $+0.0067$ & $0.500$ & $1.146$ & $0.268$ \\
defocus\_blur & 5 & $+0.0000$ & $0.229$ & $0.763$ & $0.497$ \\
elastic\_transform & 3 & $+0.0000$ & $0.706$ & $1.355$ & $0.112$ \\
elastic\_transform & 5 & $+0.0000$ & $0.264$ & $1.422$ & $0.164$ \\
fog & 3 & $+0.0133$ & $0.629$ & $1.345$ & $0.143$ \\
fog & 5 & $+0.0000$ & $0.360$ & $1.157$ & $0.264$ \\
frost & 3 & $+0.0000$ & $0.441$ & $1.187$ & $0.247$ \\
frost & 5 & $+0.0000$ & $0.331$ & $1.095$ & $0.305$ \\
gaussian\_noise & 3 & $+0.0067$ & $0.509$ & $1.266$ & $0.201$ \\
gaussian\_noise & 5 & $+0.0000$ & $0.096$ & $0.636$ & $0.545$ \\
glass\_blur & 3 & $+0.0000$ & $0.278$ & $1.076$ & $0.326$ \\
glass\_blur & 5 & $+0.0000$ & $0.141$ & $0.762$ & $0.500$ \\
impulse\_noise & 3 & $+0.0133$ & $0.502$ & $1.309$ & $0.188$ \\
impulse\_noise & 5 & $+0.0000$ & $0.113$ & $0.674$ & $0.523$ \\
jpeg\_compression & 3 & $+0.0200$ & $0.772$ & $1.372$ & $0.085$ \\
jpeg\_compression & 5 & $+0.0067$ & $0.522$ & $1.309$ & $0.184$ \\
motion\_blur & 3 & $+0.0067$ & $0.512$ & $1.206$ & $0.236$ \\
motion\_blur & 5 & $+0.0000$ & $0.220$ & $0.885$ & $0.425$ \\
pixelate & 3 & $+0.0067$ & $0.685$ & $1.378$ & $0.112$ \\
pixelate & 5 & $+0.0000$ & $0.371$ & $1.190$ & $0.258$ \\
shot\_noise & 3 & $+0.0000$ & $0.470$ & $1.244$ & $0.213$ \\
shot\_noise & 5 & $+0.0000$ & $0.113$ & $0.714$ & $0.492$ \\
snow & 3 & $+0.0000$ & $0.494$ & $1.297$ & $0.184$ \\
snow & 5 & $+0.0000$ & $0.262$ & $1.002$ & $0.349$ \\
zoom\_blur & 3 & $+0.0000$ & $0.528$ & $1.280$ & $0.190$ \\
zoom\_blur & 5 & $+0.0000$ & $0.367$ & $1.152$ & $0.272$ \\
\bottomrule
\end{tabular}
\caption{ImageNet-C, per corruption and severity, three splits. $\Delta$ worst is the change
in worst-class coverage on held-out classes at matched set size. Pooled over all $90$
conditions the improvement is $+0.003$ $[+0.001,+0.005]$.}
\label{tab:imagenetc}
\end{table}

\section{Which classes the correction moves}
\label{app:whohelped}

Worst-class coverage is a minimum, so an improvement in it says nothing on its own about the
rest of the label space. \Cref{tab:whohelped} therefore reports, for the primary ImageNet
configuration, the mean change in per-class coverage grouped into prevalence quintiles, along
with the fraction of classes whose coverage rises and falls. Numbers are from the first split,
which is the only one for which we stored the per-class breakdown.

The pattern is the one an equity intervention at a fixed resource budget should produce. The
worst class rises from $0.569$ to $0.647$ and macro-coverage moves from $0.898$ to $0.896$,
while $12\%$ of classes gain coverage and $27\%$ lose a small amount of it. Since the
comparison holds average set size fixed, coverage given to the classes that were failing has to
come from somewhere, and it comes in thin slices from the classes that had coverage to spare.
Whether that trade is worth making is a decision about the application rather than about the
method, which is why we report both the minimum and the average throughout.

Prevalence carries little information in this particular table, because ImageNet's held-out
classes have between $1147$ and $1182$ training images and are therefore close to balanced.
The quintiles are included so that the same table can be read for the long-tailed datasets,
where prevalence does vary over three orders of magnitude.

\begin{table}[h]
\centering
\small
\setlength{\tabcolsep}{4pt}
\begin{tabular}{lccc}
\toprule
Prevalence quintile & Classes & Median prev. & Mean $\Delta\hat c_y$ \\
\midrule
1 (most frequent) & $60$ & $1147$ & $-0.0005$ \\
2                 & $60$ & $1163$ & $-0.0033$ \\
3                 & $60$ & $1169$ & $-0.0033$ \\
4                 & $60$ & $1174$ & $+0.0007$ \\
5 (least frequent)& $60$ & $1182$ & $-0.0020$ \\
\midrule
\multicolumn{4}{l}{Classes improved $12.3\%$, worsened $26.7\%$} \\
\multicolumn{4}{l}{Worst class $0.569 \rightarrow 0.647$, same class before and after} \\
\bottomrule
\end{tabular}
\caption{Change in per-class coverage on the $300$ held-out ImageNet classes, by prevalence
quintile, at matched average set size. The correction concentrates its effect on the failing
classes and pays for it with small reductions spread across the rest.}
\label{tab:whohelped}
\end{table}

\section{Score-level interventions do not equalise class coverage}
\label{app:score-level}

\Cref{sec:setup} argues that improving $s(x,y)$ cannot equalise coverage across classes,
because any such change moves every class in the same direction and supplies no per-class
degree of freedom. \Cref{tab:scorelevel} is the measurement behind that argument. We built four
score functions from the same backbone, the plain softmax and three aggregations over augmented
copies of the input, and evaluated each with the marginal threshold across two miscoverage
levels, three descriptor families and three splits, for $48$ configurations.

The four scores are close to indistinguishable in worst-class coverage. At $\alpha=0.10$ and
matched to within $0.03$ in average set size they give $0.534$, $0.485$, $0.534$ and $0.518$,
so the spread across four quite different score constructions is $0.049$, and two of the three
augmentation-based scores are \emph{worse} than the plain softmax. Their macro-coverage and
marginal coverage agree to three decimal places. This is what a score-level intervention looks
like when the objective is a minimum over classes.

Two honest caveats. The comparison is on one backbone and one dataset, so it supports the
structural argument and does not substitute for it. And this configuration has a median of $25$
evaluation examples per class, which \cref{sec:depth} places below the measurability threshold,
so the near-zero effect of \textsc{Pcc} in the last column of \cref{tab:scorelevel} is
uninformative about \textsc{Pcc}. That configuration is in fact the fifth instance of the
evaluation-depth pattern we found, and we found it before we understood the pattern, which is
part of why the diagnosis in \cref{sec:depth} took three attempts.

\begin{table}[h]
\centering
\small
\setlength{\tabcolsep}{3pt}
\begin{tabular}{lcccc}
\toprule
Score & Worst-class & Set size & Macro & $\Delta$ \textsc{Pcc} \\
\midrule
Softmax (LAC)          & $0.534$ & $1.256$ & $0.900$ & $-0.005$ \\
Average over augment.  & $0.485$ & $1.228$ & $0.898$ & $+0.002$ \\
Learned aggregation    & $0.534$ & $1.233$ & $0.899$ & $+0.001$ \\
Uncertainty-modulated  & $0.518$ & $1.230$ & $0.900$ & $-0.003$ \\
\bottomrule
\end{tabular}
\caption{Four score functions on the same backbone under the marginal threshold, held-out
classes, $\alpha=0.10$, three splits. Changing the score leaves worst-class coverage within
$0.049$ and does not order the scores the way marginal coverage would. The last column reports
what \textsc{Pcc} adds on top of each, in a configuration whose $25$ evaluation examples per
class put it below the depth at which the statistic can move (\cref{sec:depth}).}
\label{tab:scorelevel}
\end{table}

\section{Splitting the calibration rows}
\label{app:split}

\Cref{prop:marginal} requires the rows used to build $\tilde\delta$ and the rows used to
calibrate $c$ to be disjoint, and the configuration reported throughout \cref{sec:exp} does not
provide that: $g_\theta$, $\lambda$, $n_{\mathrm{cal}}$ and $c$ all read the same calibration
rows. We measured the cost of the reuse in \cref{app:invisible}, where \textsc{Pcc} lands
$0.0024$ of marginal coverage below the marginal threshold on the same rows.

\textsc{Pcc}-split is the variant the proposition covers. The calibration rows of each labelled
class are divided in two by a per-class random split; the first part estimates $\delta_y$, fits
$g_\theta$ and selects $\lambda$ and $n_{\mathrm{cal}}$, and the second part is used for nothing
except the conformal quantile that gives $c$. Held-out classes and evaluation rows are untouched
in both variants, so the comparison isolates the effect of the reuse. The cost is statistical:
each half has half the rows, which makes $g_\theta$ noisier and $c$ less precise, and
\cref{tab:depth} shows that halving calibration depth from $100$ to $50$ examples per class
costs about $0.019$ of worst-class improvement.

\TODO{Fill from the PCC-split run: marginal coverage at the emitted thresholds against the
$1-\alpha$ target, the width of the finite-sample interval, and $\Delta$ worst-class on
held-out classes beside the reuse variant, ten splits. This is the only number in the paper
that is not yet measured; the run is CPU-only.}

\section{Descriptor choices}
\label{app:descriptors}

The nine features are $\|w_y\|_2$, $b_y$, the cosine distance from $w_y$ to its nearest
competing row, the mean cosine distance to its $k$ nearest rows for $k\in\{5,10,50\}$, the
spread of those distances, the cosine distance to the mean row, and $\log$ prevalence. A tenth
candidate, the cosine distance to the single nearest row computed at $k=1$, is the feature
closest in construction to the regression target and is held out of the feature set for that
reason.

\Cref{tab:ablation} ablates by feature group. Distances alone and prevalence alone are both
significantly worse than the full set, while removing prevalence is not, so the geometric
features carry the signal and the one label-derived feature is dispensable. The neighbour
counts $k\in\{5,10,50\}$ were fixed in advance and never tuned; $k=1$ enters only through the
nearest-row distance and the margin. The cross-validated error of $g_\theta$ is stable across
splits, at $0.0356$ to $0.0372$ over the five ImageNet splits, and the selected $\lambda$ is
$0.1$ in all five, which is the sense in which the fit is not delicate.

\TODO{Fill from the descriptor-sensitivity run: Euclidean in place of cosine, the head bias
dropped, and neighbour counts $\{1,2,5\}$ and $\{10,20,100\}$ in place of $\{5,10,50\}$, each at
five splits. Report $\Delta$ worst-class and the ridge coefficients so a reader can see which
features the fit actually uses.}

\section{Runtime}
\label{app:runtime}

We report the cost of a whole configuration rather than of individual operations, because
that is what the run logs measure and it is the number a practitioner cares about.
\Cref{tab:runtime} gives the median wall-clock time per configuration for each phase of the main
experiment. \textsc{Pcc} itself adds no training: the descriptors are read from the classifier's
weight matrix, $g_\theta$ is a closed-form ridge fit over at most $K$ points with $p=9$
features, and the shrinkage grid evaluates eight thresholds. No images, gradients or GPU are
involved at any point after the score dump exists.

The competitors dominate the cost. A configuration that additionally runs the seven competitor
entries takes $476$\,s against $128$\,s without them, so the comparison costs roughly $350$\,s
per split, or about $2.7$ times the whole rest of the pipeline. Most of that is the
kernel-weighted variants, which are evaluated at five bandwidths each. The complete experiment
in this paper is $228$ configurations and $5.6$ hours of CPU time.

\begin{table}[h]
\centering
\small
\setlength{\tabcolsep}{4pt}
\begin{tabular}{llrr}
\toprule
Phase & Content & Configs & Median \\
\midrule
C & competitor comparison    & $3$  & $481$\,s \\
D & ImageNet-C               & $90$ & $21$\,s \\
E & additional backbones     & $30$ & $21$\,s \\
F & evaluation-depth sweep   & $25$ & $235$\,s \\
G & long-tailed datasets     & $20$ & $63$\,s \\
H & held-out fraction sweep  & $25$ & $131$\,s \\
S & score and $\alpha$ sweep & $35$ & $127$\,s \\
\midrule
\multicolumn{3}{l}{Total, $228$ configurations} & $5.6$\,h \\
\bottomrule
\end{tabular}
\caption{Wall-clock cost per configuration, one CPU core, measured from the run reports. Phase D
and E are fast because their evaluation slices are small; phase F is slow because each of its
five depth caps is run at ten splits.}
\label{tab:runtime}
\end{table}

\section{Reproducibility}
\label{app:repro}

All $228$ configurations write one JSON report each, containing the resolved configuration,
the seed, every metric in \cref{tab:metrics} for both tables, the selected $\lambda$ and its
whole curve, the measurability regime, wall-clock time, the interpreter and package versions,
and the git commit of the code that produced it. The pre-registration documents summarised in
\cref{app:prereg} are in the same repository and are dated before the runs.

Two limitations of the comparison should be stated rather than implied. First, competitors are
invoked from their authors' released implementations, which we consider the only defensible way
to report a competitor's behaviour at $n_y=0$, but we have \emph{not} reproduced their published
numbers on their own benchmarks; our claim is about what their code does on our data, not about
whether we have implemented their method as they would. Second, \textsc{Tacp} \cite{liu2026tail}
is discussed in \cref{sec:related} and absent from \cref{tab:competitors} because its authors
release no code and it is evaluated on benchmarks we do not have, so any row we produced would
be our own reimplementation held to a lower standard than the rest of the table.
